\section{Institutional comparison: what can Colombia learn from its neighbours?}
The original benchmark uses the DPI Map \citep{ucldpimap}, operator and policy documents, expert interviews and Colombia's own findings. We keep the same five cases and the distinction between domestic system design and readiness to serve mobile or transnational users. The technical supplement records the evidence needed for the second question. A missing or unverified feature is not a low score.

Brazil's Pix illustrates a public operator's ability to set common rules and operate settlement and an alias directory \citep{bcbpix}. Peru offers a different example: the central bank phased interoperability in, starting by connecting Yape and Plin rather than asking users to move to a single national wallet \citep{bcrp2025}. Costa Rica's SINPE Móvil illustrates a simple phone-based transfer through participating institutions, but the number is linked to a local colón-denominated account \citep{bccrsinpe}. Mexico distinguishes SPEI interbank transfers from CoDi's mobile payment and collection interface \citep{banxicocodi,banxicospei}. Access to the underlying rail and use of a particular interface are different measures.

Colombia's Bre-B makes domestic interoperation available across participating providers. The comparisons help frame choices about operator responsibilities, phased connections, familiar interfaces and account requirements. They do not rank the five countries by migrant inclusion, verify acceptance of specific documents, or establish the effectiveness of cross-border protection. Any comparison of actual uptake, fees, recovery or international delivery needs its own population, period and source.
